% sesmesh — Solvent Excluded Surface meshes for molecules
% Copyright (c) 2026 Ugo Varetto
% SPDX-License-Identifier: MIT OR Apache-2.0

\section{Validation}
\label{sec:validation}

Unless stated otherwise, inputs use the default selection (first model, no
hydrogens or waters, highest-occupancy alternate locations), Bondi radii,
$r_p = 1.4$\Ang, the SES with cavities kept, the default brick size and the
automatically selected SIMD level. Tables and plots in this section are
generated from the data files in the repository by a script that ships with
the paper (\cref{app:reproduce}).

\subsection{Analytic references}
\label{sec:analytic}

A single sphere is its own SES, with area $4\pi r^2$ and volume
$\tfrac43\pi r^3$. Two spheres give an axially symmetric SES: two spherical
zones joined by the arc of the probe circle swept around the axis. Its area
$\int 2\pi R(x)\sqrt{1 + R'(x)^2}\,dx$ and volume $\int \pi R(x)^2\,dx$ are
integrated by composite Gauss--Legendre quadrature to machine precision,
split at the points where the profile is clipped by the axis in the
self-intersecting regime (the probe circle crosses the axis and the SES
splits into two components meeting at cusps); for equal radii the result
matches the closed form of the saddle area. Four configurations cover
overlapping spheres, a probe-bridged pair and a self-intersecting torus
(\cref{tab:analytic}).

\begin{table}[t]
  \centering
  \caption{Analytic references meshed at $h = 0.25$ and 0.125\Ang{}
    ($r_p = 1.4$\Ang; radii $r$ and centre distance $d$ in \AA). $C$ is the
    number of components; $A$ and $V$ the exact area and volume; relative
    errors of the mesh in \%; $p$ the observed order
    $\log_2(e_{0.25}/e_{0.125})$. Generated from \code{data/analytic.csv}.}
  \label{tab:analytic}
  \small
  \setlength{\tabcolsep}{4pt}
  \begin{tabular}{@{}lrrrrrrrrr@{}}
    \toprule
    & & & & \multicolumn{3}{c}{Area error} & \multicolumn{3}{c}{Volume error} \\
    \cmidrule(lr){5-7}\cmidrule(l){8-10}
    Configuration & $C$ & $A$ (\AA$^2$) & $V$ (\AA$^3$) &
      0.25\Ang & 0.125\Ang & $p$ & 0.25\Ang & 0.125\Ang & $p$ \\
    \midrule
    \tablerows{generated/tab-analytic.tex}
    \bottomrule
  \end{tabular}
\end{table}

The mesh is \emph{inscribed}: its vertices lie on the surface and its flat
triangles cut the convex parts, so it underestimates both area and volume
by $O(h^2)$ (\cref{fig:convergence}a). At $h = 0.25$\Ang{} the area error is
$-$\val{an-area-lo-0.25} to $-$\val{an-area-hi-0.25}\,\% and the volume
error $-$\val{an-vol-lo-0.25} to $-$\val{an-vol-hi-0.25}\,\%; at
0.125\Ang{} they fall to $-$\val{an-area-lo-0.125} to
$-$\val{an-area-hi-0.125}\,\% and $-$\val{an-vol-lo-0.125} to
$-$\val{an-vol-hi-0.125}\,\%, with observed orders between
\val{an-order-min} and \val{an-order-max}. Small atom-sized spheres are the
worst case for an inscribed mesh; the specification's accuracy requirement
for these references is therefore an $O(h^2)$ bound of 0.7\,\% for area and
1.5\,\% for volume at 0.25\Ang, which every case meets. Proteins behave
much better (\cref{sec:proteins}).

\paragraph{Vertices and normals.}
Over all meshes of \cref{tab:analytic} at the four spacings
$h \in \{1, 0.5, 0.25, 0.125\}$\Ang, the largest field value at a vertex,
$\max|F(\vect v)|$, i.e.\ the distance to the exact SES from inside
(\ref{p:dist}), is $\val{an-f-over-h}\,h$: the edge clamp of
\cref{sec:refine}, not the refinement tolerance. The largest angle between
a mesh normal and the analytic normal at the vertex position is
\val{an-normal-max}$^\circ$.

\begin{figure}[t]
  \centering
  \begin{tikzpicture}
    \begin{groupplot}[group style={group size=3 by 1, horizontal sep=19mm},
        width=0.32\linewidth, height=0.32\linewidth,
        xmode=log, log basis x=2, xlabel={$h$ (\AA)},
        xtick={0.125,0.25,0.5,1}, xticklabels={0.125,0.25,0.5,1},
        grid=major, grid style={black!10}]
      % (a) analytic references
      \nextgroupplot[ymode=log, log ticks with fixed point, ylabel={$|$relative error$|$ (\%)},
        title={\small (a) analytic references},
        ymin=0.05, ymax=30, legend pos=south east, legend cell align=left]
      \foreach \c in {0,...,4} {
        \addplot[mark=*, mark size=1.2pt, blue!70!black, thin, forget plot]
          table[x=h, y=a\c] {generated/analytic-conv.dat};
        \addplot[mark=square*, mark size=1.2pt, orange!80!black, thin, forget plot]
          table[x=h, y=v\c] {generated/analytic-conv.dat};
      }
      \addlegendimage{mark=*, blue!70!black}\addlegendentry{area}
      \addlegendimage{mark=square*, orange!80!black}\addlegendentry{volume}
      \addplot[dashed, black, domain=0.11:1.1, samples=2] {8*x^2};
      \addlegendentry{$\propto h^2$}
      % (b) proteins: area
      \nextgroupplot[ymode=log, log ticks with fixed point, ylabel={$|$area error$|$ (\%)},
        title={\small (b) \val{corpus-n} structures, area},
        ymin=0.1, ymax=5, legend pos=south east, legend cell align=left]
      \addplot[only marks, mark=o, mark size=1pt, blue!60!black, opacity=0.5]
        table[x=h, y=area] {generated/conv-entries.dat};
      \addlegendentry{entries}
      \addplot[dashed, black, domain=0.22:1.1, samples=2] {3.2*x^2};
      \addlegendentry{$\propto h^2$}
      % (c) proteins: signed volume error
      \nextgroupplot[ylabel={signed volume error (\%)},
        title={\small (c) \val{corpus-n} structures, volume},
        ymin=-1.5, ymax=0.3, legend pos=south east, legend cell align=left]
      \addplot[name path=lo, draw=none, forget plot]
        table[x=h, y=vol_min] {generated/conv-summary.dat};
      \addplot[name path=hi, draw=none, forget plot]
        table[x=h, y=vol_max] {generated/conv-summary.dat};
      \addplot[orange!25, forget plot] fill between[of=lo and hi];
      \addplot[mark=square*, mark size=1.2pt, orange!80!black]
        table[x=h, y=vol_mean] {generated/conv-summary.dat};
      \addlegendentry{mean}
      \addlegendimage{area legend, fill=orange!25, draw=none}
      \addlegendentry{range}
    \end{groupplot}
  \end{tikzpicture}
  \caption{Convergence. (a)~Absolute relative error of area (circles) and
    volume (squares) for the five analytic references of
    \cref{tab:analytic}. (b)~Area error of each of the \val{corpus-n}
    validation structures relative to its $h \to 0$ extrapolation.
    (c)~Signed volume error: mean and range over the corpus. The
    extrapolation is fitted on $h = 0.5$, 0.35 and 0.25\Ang, so the points
    at 1 and 0.7\Ang{} are out-of-sample.}
  \label{fig:convergence}
\end{figure}

\subsection{Proteins and nucleic acids}
\label{sec:proteins}

No closed form exists for real molecules, and no external surface program
was available when this study was made (\cref{sec:comparison}). Because the
vertices of \sesmesh{} lie on the exact SES, the area and volume of its
inscribed mesh converge as $A(h) = A_0 + c\,h^2$; we take the least-squares
fits $A_0$ and $V_0$ over $h \in \{0.5, 0.35, 0.25\}$\Ang{} as references
and report errors relative to them. This measures the \emph{discretization}
error of the mesh. It cannot detect a difference in the \emph{definition}
of the surface (radii, selection, probe-region details) with respect to
other programs; that requires the cross-program comparison of
\cref{sec:comparison}.

The validation corpus has \val{corpus-n} PDB entries from
\val{corpus-atoms-min} atoms (\val{corpus-atoms-min-id}) to
\val{corpus-atoms-max} atoms (\val{corpus-atoms-max-id}): peptides, small
and medium proteins, haemoglobins, nucleic acids, nucleosomes, antibodies,
a spike trimer, GroEL--GroES, ribosomal subunits and the human 80S ribosome
(\cref{app:corpus}). \Cref{tab:accuracy} and \cref{fig:convergence}b,c
summarize the errors. The area is always underestimated, and halving $h$
from 0.5 to 0.25\Ang{} divides the mean area error by \val{area-ratio-half};
the slope of the mean area error over all five spacings is
\val{area-order}. The volume error is an order of magnitude smaller than
for single spheres because convex and concave patches bias it in opposite
directions: inscribed triangles lie inside the solid on convex patches and
on the solvent side on concave ones; small, highly curved molecules
(peptides and the
\pdb{1BNA} DNA dodecamer) have the largest volume errors. The
specification's accuracy targets for proteins---area within 1\,\% and
volume within 0.5\,\% at $h = 0.25$\Ang, 3\,\% and 1\,\% at 0.5\Ang---are
met with a wide margin with respect to this discretization reference.

\begin{table}[t]
  \centering
  \caption{Mesh area and volume error relative to the $h \to 0$
    extrapolation over the \val{corpus-n} validation structures (in \%;
    ``worst'' is the most negative entry, ``max'' the most positive).
    Generated from \code{docs/validation/accuracy.csv}.}
  \label{tab:accuracy}
  \small
  \begin{tabular}{@{}rrrrrr@{}}
    \toprule
    & \multicolumn{2}{c}{Area error} & \multicolumn{3}{c}{Volume error} \\
    \cmidrule(lr){2-3}\cmidrule(l){4-6}
    $h$ (\AA) & mean & worst & mean & worst & max \\
    \midrule
    \tablerows{generated/tab-accuracy.tex}
    \bottomrule
  \end{tabular}
\end{table}

\subsection{Vertex density}
\label{sec:density}

The level-of-detail model of \cref{sec:lod} predicts
$\text{vertices}\cdot h^2/\text{area} = 1.5$ for Marching Cubes on a smooth
surface of random orientation. Over the corpus the measured constant is
\val{dens-pooled-mean}~$\pm$~\val{dens-pooled-sd} (mean $\pm$ standard
deviation for $h \le 0.5$\Ang; \cref{tab:density}), so the density mode
uses 1.5 without calibration, and its test was tightened to hit the
requested density within $\pm5$\,\%.

\begin{table}[t]
  \centering
  \caption{Vertex density constant $\text{vertices}\cdot h^2/\text{area}$
    over the \val{corpus-n} validation structures (theory: 1.5).}
  \label{tab:density}
  \small
  \begin{tabular}{@{}rrrr@{}}
    \toprule
    $h$ (\AA) & mean & SD & range \\
    \midrule
    \tablerows{generated/tab-density.tex}
    \bottomrule
  \end{tabular}
\end{table}

\subsection{Robustness}
\label{sec:robustness}

\paragraph{Random PDB entries.}
\val{sweep-n} entries were drawn at random (seed 2026) from the 263\,196
entries of the PDB at the time and meshed from mmCIF at $h = 0.5$\Ang. The
plan called for 1\,000 entries; the sweep was reduced to \val{sweep-n} by
decision. All \val{sweep-ok} produced a mesh that passed the built-in
validation (closed, consistently oriented, edge-manifold, no degenerate or
duplicate triangles), with no input errors. The entries range from
\val{sweep-atoms-min} to \val{sweep-atoms-max} atoms (median
\val{sweep-atoms-median}); together they have \val{sweep-cavities}
cavities, and one entry has \val{sweep-components-max} components. The
whole sweep took \val{sweep-seconds-total}\,s of wall time; the slowest
entry, \pdb{\val{sweep-slowest-id}} (\val{sweep-slowest-atoms} atoms,
\val{sweep-slowest-tri}\,M triangles), took \val{sweep-seconds-max}\,s. The
component counts include the sub-grid fragments of \cref{sec:fragments}
(1\,026 in 119 entries, at most 2.7\AA$^2$ each).

\paragraph{Degenerate geometry.}
Synthetic inputs that are known to be hard for analytic methods are part of
the test suite: cubic lattices (many exactly equal distances and four-sphere
coincidences), symmetric cages made of the vertices of an icosahedron and a
dodecahedron, collinear chains, exact duplicates and spheres nested inside
others, and random protein-like chains, each for the SES, SAS and vdW
surfaces. Further tests place a molecule $10^4$\Ang{} from the origin (the
single-precision kernel frame follows the molecule), use probes of 0.1 and
10\Ang, and use a 5\,000\Ang{} extent, which selects the double-precision
sampling path automatically. The C$_{60}$ fullerene of NanoShaper's examples, on which MSMS misses the
internal cavity~\cite{decherchi2013}, is part of the comparison's synthetic
set (\cref{sec:cmp-cavities}): \sesmesh{} meshes it with its cavity,
like NanoShaper and ESES.
